% !PRINTED 197
\weberpage{197}
%<<CONT
ist. Es ist leicht einzusehen, dass dieser Fall, wenn $\Theta$ einen
endlichen Grad hat, nicht vorkommen kann.

Denn ist zunächst $\beta$ oder $\gamma = 0$, so folgt aus (7), dass
$\alpha = \delta$ sein muss, und beide können $= 1$ angenommen werden.

Wenn nun $\Theta$ eine der beiden Substitutionen
\[
\begin{pmatrix} 1, & 0 \\ \gamma, & 1 \end{pmatrix}, \quad
\begin{pmatrix} 1, & \beta \\ 0, & 1 \end{pmatrix}
\]
ist, so ist für jedes $\lambda$:
\[
\Theta^\lambda = \begin{pmatrix} 1, & 0 \\ \lambda\gamma, & 1 \end{pmatrix} \text{ oder } \begin{pmatrix} 1, & \lambda\beta \\ 0, & 1 \end{pmatrix},
\]
wie man leicht durch vollständige Induction findet. Es kann
also, da $\beta$ und $\gamma$ nicht zugleich Null sein können, wenn $\Theta$ nicht
die identische Substitution ist, $\Theta^\lambda$ für kein positives $\lambda$ gleich 1
werden, wie es doch sein müsste, wenn $\lambda$ gleich dem Grade von
$\Theta$ wäre.

Ist aber $\beta$ von Null verschieden, so setzen wir, indem wir
jetzt $\Delta = 1$ und nach (8) $\alpha + \delta = 2$ annehmen,
\[
L = \begin{pmatrix} \beta, & 0 \\ 1-\alpha, & \beta^{-1} \end{pmatrix},
\]
und erhalten
\[
\begin{aligned}
L^{-1}\Theta L &= \begin{pmatrix} \beta^{-1}, & 0 \\ \alpha-1, & \beta \end{pmatrix} \begin{pmatrix} \alpha, & \beta \\ \gamma, & \delta \end{pmatrix} \begin{pmatrix} \beta, & 0 \\ 1-\alpha, & \beta^{-1} \end{pmatrix} \\
&= \begin{pmatrix} 1, & \beta^{-1} \\ 0, & 1 \end{pmatrix},
\end{aligned}
\]
also
\[
\Theta = L \begin{pmatrix} 1, & \beta^{-1} \\ 0, & 1 \end{pmatrix} L^{-1},
\]
und daraus
\[
\Theta^\lambda = L \begin{pmatrix} 1, & \lambda\beta^{-1} \\ 0, & 1 \end{pmatrix} L^{-1},
\]
was wieder für kein positives $\lambda$ gleich 1 werden kann. Wir
haben daher den Satz:

\par\noindent 1.\ \gesperrt{Jede der $n-1$ nicht identischen Substitutionen
einer Gruppe $n^{\text{ten}}$ Grades hat zwei von einander
verschiedene Pole.}

Jede nicht identische Substitution der Gruppe $G$ giebt uns
also zwei Pole. Es kann aber ein und derselbe Werth bei
mehreren verschiedenen Substitutionen als Pol auftreten. Zählen
wir einen dieser Werthe $h$ mal, wenn er in $h$ Substitutionen der
Gruppe als Pol vorkommt, so ergiebt sich die Anzahl der Pole
%CONT>>
